Let \(\phi\) be the standard Gaussian density. Choose \(K+2\) pairwise disjoint bounded intervals \(I_0,\ldots,I_{K+1}\), each of positive Gaussian probability. The \((K+1)\)-by-\((K+2)\) matrix
\[
 A_{r\ell}=\int_{I_\ell}x^r\phi(x)\,dx,
 \qquad 0\le r\le K,
\]
has a nonzero null vector \(b=(b_0,\ldots,b_{K+1})\). Rescale it so that \(\max_\ell|b_\ell|=1\), and define the bounded, nonzero function \(h=\sum_\ell b_\ell\mathbf 1_{I_\ell}\). Then
\[
 \int x^rh(x)\phi(x)\,dx=0,
 \qquad 0\le r\le K. \tag{22}
\]
For \(0<t<1\), set \(f_t(x)=\phi(x)(1+th(x))\). The case \(r=0\) of (22) makes \(f_t\) integrate to one, and \(|h|\le1\) makes it strictly positive. Since \(h\) is nonzero on a set of positive Gaussian measure, \(f_t\ne\phi\). Equation (22) says that its first \(K\) moments, including mean zero and variance one, equal the Gaussian moments. Cumulants through order \(K\) are universal polynomials in the moments through the same order, so they also agree.

Moreover \(f_t\le2\phi\), whence its even moments satisfy \(m_{2s}(f_t)\le2(2s-1)!!\). It follows that \(m_{2s}(f_t)^{-1/(2s)}\) is bounded below by a positive constant times \(s^{-1/2}\), and therefore the Carleman series diverges. The cited Carleman criterion makes \(F_t\) moment determinate. Finally,
\[
 d_{\mathrm{TV}}(F_t,\mathcal N(0,1))
 =\frac t2\int|h(x)|\phi(x)\,dx\longrightarrow0
\]
as \(t\downarrow0\). Choosing \(t\) small enough proves the claim.